\section{De la detección a la prevención: Intervention Ladder y capacidades del usuario}

La mayoría de los sistemas anteriores identifican y atienden el riesgo después de que aparece; esta sección amplía el objetivo hacia la prevention. Prevenir no significa «predecir todo lo malo», sino diseñar la friction adecuada para cada grado de certeza: educación del usuario, privacidad predeterminada, adecuación a la edad, accesos para denunciar, avisos de riesgo, restricciones de contacto, hold de contenido y escalamiento a especialistas. Cuanto más irreversible es una acción, mayores son la evidencia y la revisión que requiere.

\subsection{La intervention ladder de cuatro niveles}

Una ladder práctica puede escalar de forma gradual de educate a disrupt, protect y, por último, investigate/report. La education ayuda a los usuarios a reconocer la manipulación y los canales para pedir ayuda; la disruption eleva el costo del contacto masivo, de los mensajes de desconocidos o de la migración rápida a otra plataforma; la protection ofrece configuraciones predeterminadas, block/report y trusted-contact support; solo los eventos de alta confianza pasan a revisión especializada, a acciones de la plataforma y a los procedimientos legales.

\lecturefigure{14-prevention-ladder.png}{La prevención requiere una combinación por niveles de educación, friction de producto, escalamiento protector y respuesta legal.}{Subtítulos locales 33:10--35:30; NCMEC NetSmartz y discusión en clase; reconstrucción conceptual.}

\begin{knowledgebox}{Lectura de la figura: cuanto más débil es la evidence, más reversible debe ser la acción}
Las señales de baja confianza justifican mostrar avisos de seguridad, limitar que desconocidos encuentren al usuario u ofrecer block/report; el riesgo medio puede añadir friction, retrasar el envío o pedir una confirmación adicional; solo los eventos de alta confianza y ya revisados pasan a acciones obligatorias y a reportes. Este orden protege a los usuarios de los falsos positivos y permite que el sistema siga ofreciendo ayuda desde las etapas tempranas. El software puede reducir las oportunidades y el tiempo de respuesta, pero no puede sustituir a la familia, la escuela, las líneas de ayuda ni los servicios profesionales.
\end{knowledgebox}

\teachervoice{Cordua resume la dirección como el paso de la reactive detection a la prevention, pero no afirma que la tecnología pueda resolver el problema por sí sola. La alfabetización digital de padres y jóvenes, los canales claros para pedir ayuda y las protecciones predeterminadas de la plataforma siguen siendo parte del sistema, y no «etapas manuales que pueden eliminarse una vez que el modelo entra en producción».}

\subsection{Resumen de la sección}

Un sistema de prevención debe hacer corresponder la solidez de la evidencia con la reversibilidad de la acción. La education, la friction, el user control, la specialist review y la lawful response forman juntos la línea de defensa; ningún classifier por sí solo puede cubrir toda esta ladder.
